\documentclass[a4paper,12pt]{article}

\usepackage[utf8]{inputenc}
\usepackage[spanish, es-noshorthands]{babel}
\usepackage{geometry}
\usepackage{graphicx}
\usepackage{multirow}
\usepackage{array}
\usepackage{fancyhdr}
\usepackage{lastpage}
\usepackage{hyperref}
\usepackage{float}
\usepackage{caption}
\usepackage{booktabs}
\usepackage[table]{xcolor}
\usepackage[T1]{fontenc}
\usepackage{tabularx}
\usepackage{pdflscape}

\hypersetup{colorlinks=true,linkcolor=blue,urlcolor=blue}

\definecolor{SKTPrimary}{HTML}{3F4A4F}   % propuesta primario
\definecolor{SKTDark}{HTML}{242B2E}      % propuesta secundario
\definecolor{SKTAccent}{HTML}{979C9D}    % acento gris
\definecolor{SKTLight}{HTML}{CDCFD2}     % claro
\definecolor{SKTMid}{HTML}{6E7375}       % matiz intermedio

\geometry{
top=2.5cm,
bottom=2.5cm,
left=3cm,
right=3cm,
includehead,
headheight=130pt,
headsep=1cm
}
% ---------------- ENCABEZADO ----------------
\pagestyle{fancy}
\fancyhf{}
\renewcommand{\headrulewidth}{0pt}

\fancyhead[C]{%
\small
\setlength{\tabcolsep}{4pt}
\renewcommand{\arraystretch}{1.1}
\begin{tabular}{|m{3cm}|m{7.5cm}|m{3.8cm}|}
\hline
\multirow{5}{*}{
\parbox[c][3.5cm][c]{2.5cm}{
\centering
\includegraphics[width=2.2cm]{SKTLogo.png}
}
}
& \centering\textbf{SKT Software Solution}
& \parbox[c]{\linewidth}{\centering
\textbf{CÓDIGO:}\\
DI-GDR-PLA-001
\vspace{3pt}
} \\ \cline{2-3}

& \centering\textbf{PROCESO: Gestión de Proyectos}
& \parbox[c]{\linewidth}{\centering
\textbf{VERSIÓN:}\\
1.0
\vspace{3pt}
} \\ \cline{2-3}

& \centering Plan de Comunicación
& \parbox[c]{\linewidth}{\centering
\textbf{ELABORACIÓN:}\\
15/05/2026
} \\ \cline{2-3}

& \centering Plan y Matriz de Comunicación
& \parbox[c]{\linewidth}{\centering
\textbf{ÚLTIMA REVISIÓN:}\\
15/05/2026
} \\ \hline
\end{tabular}%
}

\begin{document}
\raggedbottom

\section{Plan de Comunicación}
Este documento establece el plan de comunicación para el desarrollo del Sistema de Gestión de Citas Médicas "Lupsi".

\section{Introducción}
El éxito del proyecto para el desarrollo del Sistema de Gestión de Citas Médicas "Lupsi" depende en gran medida de una comunicación efectiva, oportuna y transparente entre todos los interesados. En proyectos de desarrollo de software con metodologías ágiles, la sincronización entre el equipo técnico, el cliente final y los evaluadores académicos es un factor crítico de éxito. Este documento define los canales, la frecuencia y los responsables de transmitir la información clave a lo largo de todo el ciclo de vida del proyecto, asegurando el alineamiento de expectativas, la mitigación de riesgos y la resolución ágil de impedimentos técnicos o de negocio.

\section{Objetivo}
\textbf{Objetivo General:} Establecer un marco de referencia estandarizado y automatizado que garantice el flujo de información adecuado entre todos los actores del proyecto Lupsi, definiendo claramente qué se debe comunicar, a quién, cómo, cuándo y quién es el responsable de la emisión y recepción del mensaje.

\textbf{Objetivos Específicos:}
\begin{itemize}
    \item Formalizar las interacciones académicas y de evaluación con el Sponsor del proyecto.
    \item Definir los canales de demostración y retroalimentación iterativa con el cliente final.
    \item Estandarizar la comunicación interna del equipo de desarrollo, integrando herramientas de control de versiones y tableros ágiles.
    \item Automatizar la generación de reportes y alertas de bloqueo mediante la integración de un Agente virtual en plataformas de mensajería.
\end{itemize}

\section{Definiciones}
Para la correcta interpretación de este plan, se establecen los siguientes términos:
\begin{itemize}
    \item \textbf{Sponsor Académico:} Rol que provee el respaldo institucional, evalúa el cumplimiento de los hitos técnicos y aprueba formalmente las entregas.
    \item \textbf{Cliente Final:} Rol que representa a los usuarios del sistema, provee las reglas del negocio, valida el diseño y otorga retroalimentación funcional (Centro Médico Lupsi).
    \item \textbf{Sprint:} Iteración de desarrollo corta y estructurada (típicamente 1 a 2 semanas) característica del marco de trabajo ágil, al final de la cual se entrega una porción funcional de software.
    \item \textbf{Daily Sync / Stand-up:} Reunión diaria o reporte asíncrono para sincronizar al equipo técnico, donde cada desarrollador reporta su progreso y levanta la mano ante bloqueos en el código.
    \item \textbf{Agente Lupsi:} Entidad robótica automatizada integrada en Telegram que facilita el registro de las \textit{Daily Syncs} y genera documentos PDF ejecutivos sin intervención humana.
\end{itemize}

\section{Alcance}
Este plan de comunicación aplica a todas las fases del ciclo de vida del desarrollo del proyecto Lupsi, desde la toma de requerimientos iniciales hasta la fase de pruebas de estrés, despliegue y cierre formal del proyecto. Involucra a todos los miembros del equipo de desarrollo de SKT Software Solution, a la representación del Centro Médico Lupsi y al cuerpo docente evaluador.

\section{Público Objetivo}
Los principales actores involucrados en este plan son:
\begin{itemize}
    \item \textbf{Sponsor Académico:} Ing. Paulo Torres, encargado de la evaluación formal y validación de los entregables documentales.
    \item \textbf{Cliente Final:} Dra. Luz Marina Saenz, representante del Centro Médico Lupsi, principal interesada en la funcionalidad y diseño del software.
    \item \textbf{Gestor del Proyecto:} Angel Ayuquina, responsable principal de la articulación y cumplimiento de este plan.
    \item \textbf{Equipo de Desarrollo:} Sebastián Ortiz, Daniel Luisa y Alex Huachi, responsables de la ejecución técnica y el reporte continuo de avance operativo.
\end{itemize}

\section{Herramientas y Canales}
Para garantizar la fluidez de la información, el equipo ha adoptado un stack tecnológico que incluye:
\begin{itemize}
    \item \textbf{Reuniones Presenciales:} Exclusivas para entregas formales, firmas de actas y defensa académica.
    \item \textbf{Plataformas de Videoconferencia (Meet/Zoom/Teams):} Para demostraciones (Demos) quincenales y resolución de problemas técnicos complejos.
    \item \textbf{Mensajería Instantánea (WhatsApp/Telegram):} Para consultas ágiles, resolución de bloqueos diarios y recepción de notificaciones del Agente Lupsi.
    \item \textbf{Herramientas de Gestión (Trello y GitHub):} Para el seguimiento milimétrico del estado de las tareas, control de código fuente y métricas de avance del sprint.
\end{itemize}

\section{Matriz de Comunicación}
A continuación se presenta la matriz formal de comunicación, la cual detalla las estrategias adoptadas por el equipo para mantener a todos los interesados informados y coordinados.

\vspace{0.5cm}
\noindent
\makebox[\textwidth][c]{
{\footnotesize
\renewcommand{\arraystretch}{1.5}
\begin{tabular}{|>{\raggedright\arraybackslash}p{4.4cm}|>{\raggedright\arraybackslash}p{3cm}|>{\raggedright\arraybackslash}p{3.4cm}|>{\raggedright\arraybackslash}p{4cm}|>{\raggedright\arraybackslash}p{2.5cm}|}
\hline
\rowcolor{SKTPrimary}
\textcolor{white}{\textbf{QUÉ COMUNICAR}} & \textcolor{white}{\textbf{PERSONAS A QUIEN COMUNICAR}} & \textcolor{white}{\textbf{RESPONSABLE DE LA COMUNICACIÓN}} & \textcolor{white}{\textbf{CÓMO COMUNICAR (Medio/Formato)}} & \textcolor{white}{\textbf{FRECUENCIA}} \\ \hline

\textbf{1. Evaluación Académica y Aprobación Formal} \newline \textit{(Presentación de documentos, revisión del estado global y rúbricas)} & Sponsor (Docente Ing. Paulo Torres) & Gestor Angel Ayuquina & \textbf{Reunión Presencial} y entrega de actas firmadas & \textbf{Mensual} \\ \hline

\textbf{2. Seguimiento y Demostración de Avances (Demo)} \newline \textit{(Mostrar módulos funcionales, recopilar feedback y validar diseño)} & Cliente Final (Dra. Luz Marina Saenz) & Gestor Angel Ayuquina & Reunión Virtual (Meet/Zoom/Teams) & \textbf{Quincenal} \textit{(Al final de cada Sprint)} \\ \hline

\textbf{3. Consultas Funcionales y Requerimientos} \newline \textit{(Aclaración de reglas de negocio, horarios, flujos del centro médico)} & Cliente Final y Equipo de Desarrollo & Gestor Angel Ayuquina & Grupo de WhatsApp con el Cliente o Reunión Corta & Según sea necesario (Ad-hoc) \\ \hline

\textbf{4. Asignación y Seguimiento Operativo Interno} \newline \textit{(Coordinación de código, backend/frontend, actualización de tareas)} & Equipo de Desarrollo (Sebastián, Daniel, Alex) & Todos los Desarrolladores & Tablero Kanban en \textbf{Trello} y control de ramas en \textbf{GitHub} & Diaria / Continua \\ \hline

\textbf{5. Resolución de Bloqueos Internos (Daily Sync)} \newline \textit{(Standup diario y alertas de estancamiento en el código)} & Equipo de Desarrollo Lupsi y Gestor PM & \textbf{Agente Lupsi (Automático)} y Desarrolladores & Interacción con el bot en el grupo de \textbf{Telegram} \newline (Comando \texttt{/standup}) & Diaria \\ \hline

\textbf{6. Resultados de Calidad y Estado Técnico} \newline \textit{(Progreso del código, métricas, presupuesto y cumplimiento en Trello)} & Sponsor (Docente), Cliente y Equipo & \textbf{Agente Lupsi (Automático)} & \textbf{Reporte Ejecutivo en PDF} generado por el bot y enviado vía Telegram & Al finalizar cada sprint \newline (Semanal) \\ \hline

\textbf{7. Cierre Formal del Proyecto} \newline \textit{(Defensa final del software funcionando, entrega de código y manuales)} & Sponsor (Docente) y Cliente Final & Gestor Angel Ayuquina & Presentación formal presencial (Defensa) & Al finalizar el proyecto \\ \hline
\end{tabular}
}
}
\vspace{0.5cm}

\section{Firmas de Responsabilidad}

\renewcommand{\arraystretch}{2}

\noindent
\begin{tabular}{
p{0.28\textwidth}
p{0.30\textwidth}
>{\centering\arraybackslash}p{0.22\textwidth}
>{\centering\arraybackslash}p{0.12\textwidth}
}
\hline
\textbf{Rol} &
\textbf{Nombre / Representante} &
\textbf{Firma (QR)} &
\textbf{Fecha} \\
\hline

Gestor del Proyecto
& Angel Ayuquina
& \includegraphics[width=3.2cm]{QR_Sello_Angel_Ayuquina.png}
& 25/03/2026 \\[0.6cm]

Desarrollador Backend
& Sebastián Ortiz
& \includegraphics[width=3.2cm]{QR_Sello_Sebastían_Ortiz.png}
& 25/03/2026 \\[0.6cm]

Desarrollador Fullstack
& Daniel Luisa
& \includegraphics[width=3.2cm]{QR_Sello_Daniel_Luisa.png}
& 25/03/2026 \\[0.6cm]

Desarrollador Frontend
& Alex Huachi
& \includegraphics[width=3.2cm]{QR_Sello_Alex_Huachi.png}
& 25/03/2026 \\[0.6cm]

\hline
\end{tabular}

\end{document}